\subsection{Étude de cas sur les préjudices algorithmiques : le scandale des allocations familiales aux Pays-Bas}

Le scandale néerlandais des allocations familiales, connu sous le nom de «Toeslagenaffaire», constitue un exemple tragique et concret des dangers de la bureaucratie algorithmique non contrôlée. À partir de 2013, l'administration fiscale néerlandaise a utilisé un algorithme d'apprentissage automatique pour créer des profils de risque afin de détecter les fraudes aux allocations de garde d'enfants \citet{Heikkilä2023How}.

Le système signalait les demandeurs sur la base de critères qui incluaient la double nationalité et un faible revenu \citet{Heikkilä2023How}. Des dizaines de milliers de familles, de manière disproportionnée issues de minorités ethniques, ont été faussement accusées de fraude sur la base des indicateurs de risque du système. Elles ont été contraintes de rembourser des dizaines de milliers d'euros, les plongeant dans la pauvreté et causant des drames humains, y compris des suicides et le placement de plus de 1000 enfants en famille d'accueil \citep{TaxAdmin.ai2025The, AmnestyInternational2021}.

L'analyse de ce scandale révèle plusieurs défaillances systémiques. 
\begin{enumerate}
    \item D'abord, l'algorithme a institutionnalisé et amplifié un biais discriminatoire. L'utilisation de la double nationalité comme indicateur de risque reflétait et renforçait les stéréotypes raciaux \citep{Alrefai2025The}. 
    L'administration fiscale tenait également des listes noires secrètes, où la nationalité turque ou marocaine était un facteur de risque, et être sur cette liste augmentait le score de risque dans le système d'allocations \citet{Heikkilä2023How}. 
    \item  Par ailleurs, il y avait un manque flagrant de transparence et de recours. Les citoyens n'avaient aucun moyen de savoir pourquoi ils étaient signalés ou de se défendre efficacement \citet{Heikkilä2023How}. Qui plus est, les fonctionnaires chargés de l'examen des cas signalés, n'avaient pas d'informations sur les raisons pour lesquelles le système avait généré un score de risque élevé, les conduisant souvent à suivre aveuglément la recommandation de l'algorithme \citet{AmnestyInternational2021}.
\end{enumerate}

La \textbf{Toeslagenaffaire} doit nous sensibiliser sur les risques de l'automatisation sans garanties adéquates. Soulignant la nécessité d'une surveillance humaine vigoureuse, de la transparence algorithmique, de l'occasion de contester les décisions automatisées et de la proscription d'employer des \textbf{proxys discriminatoires} (p. ex. l'origine ethnique) \citep{Heikkilä2023How, AmnestyInternational2021}.

\section{La prérogative de l'état réimaginée : du monopole de la violence de Weber à la gouvernance algorithmique}
\label{sec:prerogative}

Cette section fondamentale établit le cadre théorique central du rapport: le <<~wébérisme numérique~>>. Il y est soutenu que la police prédictive n'est pas simplement un nouvel outil pour l'état, mais représente une évolution fondamentale de la nature même du pouvoir étatique, passant d'une base purement \textbf{légale-rationnelle} à une base \textbf{algorithmique-rationnelle}.

\subsection{L'état wébérien : rationalisation, bureaucratie et force légitime}
\label{ssec:etat-weberien}

La définition classique de l'état moderne, telle que formulée par Max Weber, repose sur sa revendication du <<~monopole de la violence physique légitime~>> sur un territoire donné \citep{Weber1960Le, Weber2002Economy}.

Au sein des frontières nationales, la police est l'institution principale qui incarne et exerce ce monopole \citep{Brodeur2010The}. La légitimité de cette violence n'est pas arbitraire: à l'ère moderne, elle découle d'une autorité de type légal-rationnel. Ce modèle d'autorité se fonde sur un système de lois impersonnelles et universellement appliquées, ainsi que sur des procédures bureaucratiques conçues pour être prévisibles et équitables. 

Le processus historique par lequel cette forme d'autorité a remplacé les dominations traditionnelles (basées sur la coutume) ou charismatiques (basées sur la dévotion à un individu) est ce que Weber a identifié comme la rationalisation.

La bureaucratie est le moteur organisationnel de cet état rationnel. Elle fonctionne comme une machine administrative
